% ****** Start of file apssamp.tex ******
%
%   This file is part of the APS files in the REVTeX 4.2 distribution.
%   Version 4.2a of REVTeX, December 2014
%
%   Copyright (c) 2014 The American Physical Society.
%
%   See the REVTeX 4 README file for restrictions and more information.
%
% TeX'ing this file requires that you have AMS-LaTeX 2.0 installed
% as well as the rest of the prerequisites for REVTeX 4.2
%
% See the REVTeX 4 README file
% It also requires running BibTeX. The commands are as follows:
%
%  1)  latex apssamp.tex
%  2)  bibtex apssamp
%  3)  latex apssamp.tex
%  4)  latex apssamp.tex
%
\documentclass[
reprint,
superscriptaddress,
%groupedaddress,
%unsortedaddress,
%runinaddress,
%frontmatterverbose, 
%preprint,
%preprintnumbers,
%nofootinbib,
%nobibnotes,
%bibnotes,
amsmath,amssymb,
aps,
%pra,
%prb,
prl,
%rmp,
%prstab,
%prstper,
floatfix,
]{revtex4-2}

\usepackage{graphicx}% Include figure files
\usepackage{dcolumn}% Align table columns on decimal point
\usepackage{bm}% bold math
\usepackage{physics}
\usepackage{mathrsfs}
\usepackage[colorlinks,linkcolor=blue,anchorcolor=blue,citecolor=blue,breaklinks=true]{hyperref}
\usepackage[ruled,vlined,]{algorithm2e} 
\usepackage[normalem]{ulem} % [normalem] keeps \emph as italics
\usepackage{orcidlink}

\DontPrintSemicolon

\usepackage{amsfonts,amssymb}
\usepackage{textcomp}
\usepackage{float}
\usepackage{color}
\usepackage{tikz}
\usepackage{amsthm}
\usepackage{mathtools}
\usepackage{hyperref}
\hypersetup{
    colorlinks = true,
    %linkcolor = [rgb]{0.25, 0.41, 0.88},
    %citecolor = [rgb]{0.25, 0.41, 0.88},
    urlcolor = blue
}
\urlstyle{same}
%\usepackage{caption}
%\usepackage{subcaption}
\usepackage{dsfont}
\usepackage{comment}
\newtheorem{theorem}{Theorem}[section]
\theoremstyle{definition}
\newtheorem{definition}{Definition}[section]
\newtheorem{lemma}[theorem]{Lemma}
\newtheorem{corollary}{Corollary}[theorem]
\newtheorem{proposition}[theorem]{Proposition}
\renewcommand{\qedsymbol}{$\blacksquare$}


%\usepackage{cancel}

%\numberwithin{equation}{section}

\newcommand{\vac}{\mrm{vac}}
\newcommand{\sbra}[1]{\langle #1 |}
\newcommand{\sket}[1]{| #1 \rangle}
\newcommand{\anglelr}[1]{\left\langle #1 \right\rangle}
\newcommand{\sanglelr}[1]{\langle #1 \rangle}
\newcommand{\sbraket}[2]{\sanglelr{ #1 | #2 }}
\newcommand{\dif}[1]{\frac{\mathrm{d}}{\mathrm{d}#1}}
\newcommand{\diff}[2]{\frac{\mathrm{d}#1}{\mathrm{d}#2}}
\newcommand{\gd}{\boldsymbol{\nabla}}
\newcommand{\dive}{\boldsymbol{\nabla}\cdot}
\newcommand{\lr}[1]{\left(  #1  \right)}
\newcommand{\lrm}[1]{\left[  #1  \right]}
\newcommand{\lrg}[1]{\left\{  #1  \right\}}
\newcommand{\de}{\partial}
\newcommand{\pdif}[1]{\frac{\partial}{\partial#1}}
\newcommand{\pdiff}[2]{\frac{\partial#1}{\partial#2}}
\newcommand{\mrm}[1]{\mathrm{#1}}
\newcommand{\bo}[1]{\mathbf{#1}}
\newcommand{\bog}[1]{\boldsymbol{#1}}
\newcommand{\vint}[3]{\int_{#1}#2\cdot\mathrm{d}\mathbf{#3}}
\newcommand{\voint}[3]{\oint_{#1}#2\cdot\mathrm{d}\mathbf{#3}}
\newcommand{\Vint}{\int\mrm{d}^3\bo{r}'}
\newcommand{\rr}{\mathbf{r} - \mathbf{r}'}
\newcommand{\arr}{\left|\mathbf{r} - \mathbf{r}'\right|}
\newcommand{\nx}{\hat{\bo{n}}_1}
\newcommand{\ny}{\hat{\bo{n}}_2}
\newcommand{\alert}[1]{{\color{red}#1}}
\newcommand{\valert}[1]{{\color{violet}#1}}
\newcommand{\alertV}[1]{{\color{violet}#1}}
\newcommand{\dimer}[2]{\sigma^{\alpha_{#1 #2}}_{#1}\sigma^{\alpha_{#1 #2}}_{#2}}

\newcommand{\dodecagon}{{\wordfig{0.25}{figures/dodecagon.pdf}}}

\newcommand{\wordfig}[2]{\begin{minipage}[h]{#1 cm}
	\vspace{0pt}
	\includegraphics[width=\textwidth]{#2}
   \end{minipage}}

\newcommand{\segment}[2]{\mathcal{S}_{\overrightarrow{#1#2}}}
\newcommand{\lsegment}[2]{\mathcal{S}_{\overleftarrow{#1#2}}}
    
\newcommand{\TReven}{\bo{\mathsf{Ev}}}
\newcommand{\TRodd}{\bo{\mathsf{Od}}}
\newcommand{\Good}{\bo{\mathsf{Gd}}}

\begin{document}

\preprint{APS/123-QED}

%%%%%%%%%%%%%%%%%%%%%%%%%% title %%%%%%%%%%%%%%%%%%%%%%%%%%
\title{Quantum Magic in the Kitaev Spin Liquid from the Organization of the Density Matrix}

\author{Chen-Chih Wang, 
\orcidlink{0009-0006-6656-8071}}
\affiliation{Physics Department, National Tsing Hua University, Hsinchu 30013, Taiwan}

\author{Yi-Ping Huang, \orcidlink{0000-0001-6453-1191}}
\email[Correspondence email address: ]{yphuang@phys.nthu.edu.tw} 
\affiliation{Physics Department, National Tsing Hua University, Hsinchu 30013, Taiwan}
\affiliation{Physics Division, National Center for Theoretical Sciences, Taipei 10617, Taiwan}
\affiliation{Institute of Physics, Academia Sinica, Taipei 115, Taiwan}

\author{Sungkit Yip, \orcidlink{0000-0002-3542-5186}}
\email[Correspondence email address: ]{yip@phys.sinica.edu.tw} 
\affiliation{Institute of Physics, Academia Sinica, Taipei 115, Taiwan}
\affiliation{Institute of Atomic and Molecular Sciences, Academia Sinica, Taipei 10617, Taiwan}


\date{\today}

\begin{abstract}
    
\end{abstract}

                              %display desired
%%%%%%%%%%%%%%%%%%%%%%%%%% title %%%%%%%%%%%%%%%%%%%%%%%%%%
\maketitle

%%%%%%%%%%%%%%%%%%%%%%%%% main text %%%%%%%%%%%%%%%%%%%%%%%%%

%%%%%%%%%%%%% intro %%%%%%%%%%%%%

%%%%%%%%%%%%% fundamental %%%%%%%%%%%%%
% dentiy matrix and its organization

%%%%%%%%%%%%% magic %%%%%%%%%%%%%

%%%%%%%%%%%%% result %%%%%%%%%%%%%
% numerical simulations


%%%%%%%%%%%%%%%%%%%%%%%%% main text %%%%%%%%%%%%%%%%%%%%%%%%%

%%%%%%%%%%%%%%%%%%%%%%%% bibliography %%%%%%%%%%%%%%%%%%%%%%%%
% The \nocite command causes all entries in a bibliography to be printed out
% whether or not they are actually referenced in the text. This is appropriate
% for the sample file to show the different styles of references, but authors
% most likely will not want to use it.
% \nocite{*}

% \bibliography{ref}% Produces the bibliography via BibTeX.
% \bibliographystyle{my_apsrev4-2}
%%%%%%%%%%%%%%%%%%%%%%%% bibliography %%%%%%%%%%%%%%%%%%%%%%%%

\end{document} 



